\chapter{Why dynamic pre-generation can be slower than on-demand}
\label{ch:why-dynamic}

The aim of compiler-driven pre-generation is a latency comparable to or better than on-demand generation (ODG) with fewer wasted pairs than the autonomous
protocols (CGP, ACGP). Experiment~1 found that, on the corrected scheduler and at the modelled memory counts, dynamic pre-generation is \emph{slower} than ODG on the congested benchmarks
(QFT at four memories: ODG $1.59$\,ms, dynamic $5.57$\,ms, mean of nine runs), and faster once memory is plentiful (QFT at sixteen memories: $1.14$ against $0.80$\,ms). The tables of this chapter use one mapping seed (42) unless they say otherwise, so their QFT numbers ($1.57$ and $5.46$\,ms at four memories) differ slightly. This chapter finds out why.
Everything here was measured on ACE with the corrected scheduler, and the explanation is checked against the code (Section~\ref{sec:code-reading}).

\section{What the old, serialized scheduler showed}
The scheduler used before the corrections let each core serve one request per batch. Under it the same comparison gave the opposite sign
(Table~\ref{tab:legacy-vs-corrected}): dynamic was faster than ODG at every memory count, by $41\%$ on QFT at four memories in the old scheduler's own latency.
Two things changed. First, the old latency started a request's clock when its batch was submitted, so the long waits for earlier batches were not in it; counting them
(``old common'') the gain on QFT at four memories is $11\%$, and the whole trace took $13.2$\,s with ODG and $12.3$\,s with dynamic (Table~\ref{tab:legacy-makespan}).
Second, the old scheduler let a core serve one request at a time, and in the old QFT runs at four memories no request waited for a memory under either policy
(retry wait $0.0$\,ms for ODG and for dynamic), while ODG's own latency did not respond to the number of memories at all. The corrected scheduler runs the requests of a layer
concurrently, as the hardware would, and the waiting appears.
\begin{table}[H]
\centering\small
\resizebox{\ifdim\width>\linewidth\linewidth\else\width\fi}{!}{%
\begin{tabular}{@{}llrrrrrrrrr@{}}
\toprule
Workload & Mem. & Old own ODG & Old own Dyn. & Old own gain & Old common ODG & Old common Dyn. & Old common gain & New common ODG & New common Dyn. & New common gain \\
\midrule
QFT & 2 & 1.11 & 0.93 & 17\% & 25.06 & 23.95 & 4\% & 4.72 & 10.06 & -113\% \\
QFT & 3 & 1.11 & 0.77 & 31\% & 25.06 & 23.03 & 8\% & 2.39 & 7.49 & -214\% \\
QFT & 4 & 1.11 & 0.65 & 41\% & 25.06 & 22.30 & 11\% & 1.59 & 5.57 & -250\% \\
QFT & 32 & 1.11 & 0.25 & 77\% & 25.06 & 20.28 & 19\% & 1.13 & 0.25 & 78\% \\
QAOA & 2 & 1.11 & 0.91 & 18\% & 13.39 & 12.56 & 6\% & 3.42 & 6.89 & -101\% \\
QAOA & 3 & 1.11 & 0.74 & 34\% & 13.39 & 11.84 & 12\% & 1.74 & 5.33 & -207\% \\
QAOA & 4 & 1.11 & 0.61 & 45\% & 13.39 & 11.34 & 15\% & 1.27 & 3.88 & -205\% \\
QAOA & 32 & 1.11 & 0.25 & 78\% & 13.39 & 10.14 & 24\% & 1.13 & 0.25 & 78\% \\
Draper QFT adder & 2 & 1.12 & 0.84 & 26\% & 9.40 & 8.57 & 9\% & 2.16 & 4.55 & -111\% \\
Draper QFT adder & 3 & 1.12 & 0.60 & 47\% & 9.40 & 7.91 & 16\% & 1.30 & 3.67 & -182\% \\
Draper QFT adder & 4 & 1.12 & 0.47 & 58\% & 9.40 & 7.55 & 20\% & 1.15 & 2.48 & -116\% \\
Draper QFT adder & 32 & 1.12 & 0.25 & 78\% & 9.40 & 7.05 & 25\% & 1.13 & 0.25 & 78\% \\
Graph state & 2 & 1.16 & 0.84 & 27\% & 9.98 & 9.05 & 9\% & 2.72 & 6.36 & -134\% \\
Graph state & 3 & 1.16 & 0.63 & 46\% & 9.98 & 8.42 & 16\% & 1.53 & 5.01 & -228\% \\
Graph state & 4 & 1.16 & 0.49 & 58\% & 9.98 & 7.91 & 21\% & 1.24 & 3.23 & -161\% \\
Graph state & 32 & 1.16 & 0.25 & 78\% & 9.98 & 7.28 & 27\% & 1.18 & 0.25 & 79\% \\
GHZ & 2 & 1.17 & 0.51 & 56\% & 1.17 & 0.51 & 56\% & 1.17 & 0.51 & 56\% \\
GHZ & 3 & 1.17 & 0.26 & 78\% & 1.17 & 0.26 & 78\% & 1.17 & 0.26 & 78\% \\
GHZ & 4 & 1.17 & 0.26 & 78\% & 1.17 & 0.26 & 78\% & 1.17 & 0.26 & 78\% \\
GHZ & 32 & 1.17 & 0.26 & 78\% & 1.17 & 0.26 & 78\% & 1.17 & 0.26 & 78\% \\
\bottomrule
\end{tabular}}
\caption{Request latency (ms) of on-demand (ODG) and dynamic (Dyn.) and the dynamic gain $(\text{ODG}-\text{Dyn.})/\text{ODG}$. Old own: the old serialized scheduler's own latency, whose clock for a request starts when its batch is submitted. Old common: the same old runs under the common definition, which also counts the wait for earlier batches. New common: the corrected per-qubit scheduler under the common definition. Mean of nine runs (three mapping seeds, three physics seeds).}
\label{tab:legacy-vs-corrected}
\end{table}

\begin{table}[H]
\centering\small
\resizebox{\ifdim\width>\linewidth\linewidth\else\width\fi}{!}{%
\begin{tabular}{@{}lrrrr@{}}
\toprule
Workload & Old ODG & Old Dyn. & New ODG & New Dyn. \\
\midrule
QFT & 13.23 & 12.34 & 4.66 & 5.75 \\
QAOA & 12.76 & 11.72 & 5.48 & 6.62 \\
Draper QFT adder & 7.64 & 6.97 & 4.88 & 5.03 \\
Graph state & 0.40 & 0.36 & 0.24 & 0.26 \\
GHZ & 0.33 & 0.30 & 0.33 & 0.30 \\
\bottomrule
\end{tabular}}
\caption{End-to-end time of the whole trace (s) at four memories per core, old serialized scheduler and corrected scheduler (mean of nine runs).}
\label{tab:legacy-makespan}
\end{table}


\section{What the code does with a pre-generated pair}
\label{sec:code-reading}
The explanation below was checked against the source, not inferred from the results. File and line numbers refer to the repository at the commit that contains this report.

\begin{enumerate}
  \item \textbf{When a dynamic preparation starts.} The planner gives each request a generation layer. For the dynamic policy the candidate layers are tried from the one just
    before the request's own layer, backwards (\code{compiler_trace.py:254}), and the first one with room is taken. The controller launches a layer's preparations when that layer is
    activated (\code{request_scheduler.py:151--158}, \code{compiler_scheduler.py:66}). So a pair for a request in layer $L$ is normally started at the \emph{start} of layer $L-1$
    and has to wait for all of layer $L-1$ to finish before its request can even be released. The lookahead of eight layers is only an upper limit; varying it made no difference
    (Section~\ref{sec:ablations}).
  \item \textbf{What a prepared pair occupies.} A preparation books one memory card at each of its two cores, with a window of 1\,000\,ms
    (\code{adaptive_continuous.py:192--236}), and it may not exceed the compiler cap, the same counter that limits the autonomous protocol
    (\code{adaptive_continuous.py:201}). The booking is released only when the pair is consumed (\code{reservation.py:473}) or its window ends.
    If the cap is reached the preparation is rejected and retried every millisecond until its request starts (\code{compiler_scheduler.py:113--120}).
  \item \textbf{A request that has a ready pair still needs a memory of its own.} Admission asks for a free card at both cores, and a card with any overlapping booking
    is not free (\code{request_scheduler.py:202--215, 232--255}); the prepared pair's own booking overlaps. When the request's reservation runs, the generation protocol
    looks for a ready pair for the same two cores, and if it finds one it \emph{swaps} the pair's state into the request's own memory after a classical delay
    (\code{generation.py:265--296, 301--343}) and only then releases the prepared pair's memory (\code{generation.py:342}). So for a short time a hit holds \emph{two} cards
    at each core. This is a property of ACE's design, not of the compiler policy, and it is why requests served from a ready pair can still wait for admission.
  \item \textbf{What an on-demand request occupies.} One card per core, from admission until completion. Its booking window is the request duration (50\,ms), but the booking is
    released at completion, so what matters is the time it exists, about 6\,ms in the measurements, not the window.
  \item \textbf{CGP and ACGP.} Every cycle of the autonomous protocol draws a neighbour by roulette wheel over a probability table (\code{adaptive_continuous.py:262--276}) and books one card at both
    ends until a 0.2\,s boundary within the next 0.2\,s (\code{adaptive_continuous.py:147--178}; the period is $2\times10^{11}$\,ps, \code{node.py:57}), as long as fewer than the cap are in use.
    The table is uniform for CGP (\code{adaptive_continuous.py:243--260}). For ACGP it is updated every 0.2\,s: each neighbour that appears in an entanglement path used in the last 0.2\,s gains $0.05$
    and the table is renormalized (\code{adaptive_continuous.py:404--447}). Pairs are therefore made for whichever neighbour was drawn, not for a particular request.
    An application request takes a pair for its own two cores if one exists: first a compiler pair targeted at \emph{that} request, otherwise an ordinary autonomous pair with no
    compiler metadata (\code{adaptive_continuous.py:511--523}), through the same swap. A pair nobody uses before its window ends is discarded; that is the waste the experiments count.
\end{enumerate}
Two consequences follow directly from the code and are tested in the measurements below. First, pre-generation always costs memory cards for longer than an on-demand request does
(a ready pair exists \emph{before} its request, and for a later hop of a chain well before it, because a hop is released only after its predecessor completed). Second, it
saves only the generation time (a request served from a ready pair generates in about a quarter of a millisecond instead of about a millisecond); the lead and the classical delays are the same either way.


\section{Measurements}
\label{sec:measurements}
Diagnostic runs (the code is in \code{tools/diagnostics/}; the results are listed in the repository map, Appendix~\ref{app:repo}) record, for each run, how long each request waited and why: at every
admission check that found no free memory they record which kind of booking held the memory cards at the stuck core. They use ACE as it is and change nothing in the simulation.

\subsection{The extra latency is waiting, not generation}
\begin{table}[H]
\centering\small
\resizebox{\ifdim\width>\linewidth\linewidth\else\width\fi}{!}{%
\begin{tabular}{@{}rrrrrrrrr@{}}
\toprule
Mem. & ODG lat. & ODG wait & Dynamic lat. & Dynamic wait & CGP lat. & CGP wait & ACGP lat. & ACGP wait \\
\midrule
2 & 4.60 & 3.48 & 9.43 & 8.50 & 12.29 & 11.39 & 11.86 & 10.99 \\
3 & 2.36 & 1.23 & 7.05 & 6.28 & 10.58 & 9.89 & 10.65 & 9.96 \\
4 & 1.57 & 0.44 & 5.46 & 4.81 & 10.15 & 9.58 & 10.52 & 9.98 \\
6 & 1.20 & 0.07 & 3.44 & 2.95 & 8.98 & 8.56 & 9.47 & 9.08 \\
8 & 1.15 & 0.01 & 2.62 & 2.24 & 8.92 & 8.55 & 9.26 & 8.92 \\
16 & 1.14 & 0.00 & 0.71 & 0.46 & 8.03 & 7.74 & 8.95 & 8.66 \\
\bottomrule
\end{tabular}}
\caption{QFT, mapping seed 42, mean of three physics seeds: request latency and the part of it spent waiting for a free memory (ms). The remainder of the latency is generation.}
\label{tab:diag-latency-wait}
\end{table}

Generation itself gets faster with pre-generation (QFT at four memories: $1.13$\,ms for ODG, $0.65$\,ms for dynamic on average, $0.25$\,ms for a request served from a ready pair),
but the wait for a free memory grows from $0.44$ to $4.81$\,ms. CGP and ACGP wait even more ($9.6$ and $10.0$\,ms) and improve little with more memory: their wait is
$7.7$ and $8.7$\,ms even at sixteen memories (Table~\ref{tab:diag-latency-wait}).

\subsection{What holds the memories}
\begin{table}[H]
\centering\small
\resizebox{\ifdim\width>\linewidth\linewidth\else\width\fi}{!}{%
\begin{tabular}{@{}rrrrrrrr@{}}
\toprule
Mem. & ODG & Dynamic & CGP & ACGP & Dyn. by compiler & CGP by cache & ACGP by cache \\
\midrule
2 & 18.3\% & 23.3\% & 42.2\% & 44.4\% & 40\% & 47\% & 46\% \\
3 & 12.0\% & 21.3\% & 44.8\% & 45.9\% & 55\% & 65\% & 65\% \\
4 & 6.8\% & 17.2\% & 45.3\% & 46.1\% & 60\% & 74\% & 74\% \\
6 & 1.6\% & 12.0\% & 46.5\% & 47.0\% & 66\% & 83\% & 83\% \\
8 & 0.4\% & 8.6\% & 45.9\% & 47.3\% & 73\% & 87\% & 87\% \\
16 & 0.0\% & 1.1\% & 46.5\% & 46.8\% & 81\% & 94\% & 94\% \\
\bottomrule
\end{tabular}}
\caption{QFT, mapping seed 42. Left: share of admission checks that found no free memory at one of the request's two cores. Right: of the memory cards found taken at a stuck core, the share held by the policy's own pre-generation bookings (the rest are held by other requests).}
\label{tab:diag-blocking}
\end{table}

At four memories $17\%$ of dynamic's admission checks find no free memory, against $7\%$ for ODG; of the cards taken at a stuck core, $60\%$ are held by the compiler's own prepared
pairs. For CGP and ACGP $42$ to $47\%$ of checks are stuck at every memory count, and $74\%$ of the taken cards at four memories ($94\%$ at sixteen) are autonomous bookings: the
protocol fills every memory it is allowed to, with pairs for neighbours drawn at random (Table~\ref{tab:diag-blocking}).

\subsection{Requests with a ready pair still wait}
\begin{table}[H]
\centering\small
\resizebox{\ifdim\width>\linewidth\linewidth\else\width\fi}{!}{%
\begin{tabular}{@{}lrrrrrr@{}}
\toprule
Mem. & Hits & Hit wait & Hits that waited & Other wait & Pair held & Request holds \\
\midrule
2 & 21.0\% & 4.62 & 32\% & 9.54 & 28.9 & 6.2 \\
3 & 39.0\% & 3.68 & 37\% & 7.95 & 25.9 & 6.0 \\
4 & 52.8\% & 2.93 & 35\% & 6.91 & 23.7 & 5.9 \\
6 & 73.1\% & 2.03 & 31\% & 5.48 & 21.2 & 5.7 \\
8 & 84.7\% & 1.67 & 25\% & 5.42 & 20.2 & 5.6 \\
16 & 99.1\% & 0.39 & 6\% & 7.56 & 17.2 & 5.5 \\
\bottomrule
\end{tabular}}
\caption{Dynamic, QFT, mapping seed 42 (mean of three physics seeds). Hits: share of requests served by a ready pair, i.e. a prepared pair generated by the time it was used. Hit wait and Other wait: mean wait for a free memory (ms) of the hits and of the other requests. Pair held: time (ms) a ready pair waits in memory before it is used. Request holds: time (ms) an ordinary request keeps its memory card booked.}
\label{tab:diag-hits}
\end{table}

Of the requests dynamic serves from a ready pair, $35\%$ still wait for admission at four memories (Table~\ref{tab:diag-hits}). That follows from the code: the request needs a free card of its own
in addition to the card the ready pair occupies (item~3 of Section~\ref{sec:code-reading}).
A ready pair waits $24$\,ms in memory on average before it is used at four memories; an ordinary request keeps its card booked for $5.9$\,ms.

\subsection{Why a prepared pair waits so long}
\begin{table}[H]
\centering\small
\resizebox{\ifdim\width>\linewidth\linewidth\else\width\fi}{!}{%
\begin{tabular}{@{}lrrrr@{}}
\toprule
Position in the chain & Pairs used & Made before its layer started (ms) & Ready pair held (ms) & Request eligible after pair ready (ms) \\
\midrule
first hop & 604 & 4.4 & 11.9 & 4.4 \\
second hop & 257 & 4.6 & 23.7 & 12.2 \\
third hop or later & 464 & 4.9 & 39.2 & 31.4 \\
\bottomrule
\end{tabular}}
\caption{Dynamic, QFT, four memories, mapping seed 42 (mean of three physics seeds): the pair of a later hop is made as early as the pair of a first hop, but its request becomes eligible much later, so the pair waits longer.}
\label{tab:diag-chain}
\end{table}

A pair is made at the start of the layer before its request's layer (Section~\ref{sec:code-reading}, item~1), and the request is released only when its layer has started
and, for the second and later hop of a qubit's chain, when the previous hop has completed. Table~\ref{tab:diag-chain} shows both effects: the pair of a later hop is made about as early as the pair of a first hop
but its request is eligible much later, so the pair holds a memory for $24$ and $39$\,ms instead of $12$\,ms. More than half of dynamic's pairs on QFT ($54\%$) are for hops after the first.

\section{What does and does not change the result}
\label{sec:ablations}
\begin{table}[H]
\centering\small
\resizebox{\ifdim\width>\linewidth\linewidth\else\width\fi}{!}{%
\begin{tabular}{@{}llrrrr@{}}
\toprule
Policy & Change from the baseline & Latency (ms) & Wait (ms) & Ready-pair hits & Fidelity \\
\midrule
Dynamic & baseline & 5.46 & 4.81 & 53\% & 0.9274 \\
Dynamic & compiler memories 1 of 4 & 1.88 & 0.95 & 21\% & 0.9431 \\
Dynamic & compiler memories 2 of 4 & 2.77 & 2.01 & 39\% & 0.9364 \\
Dynamic & request-duration-ms=10 & 5.45 & 4.80 & 53\% & 0.9274 \\
Dynamic & request-duration-ms=100 & 5.46 & 4.80 & 53\% & 0.9274 \\
Dynamic & request-duration-ms=20 & 5.46 & 4.80 & 53\% & 0.9274 \\
Dynamic & pairs only for the first 1 & 1.89 & 0.98 & 24\% & 0.9449 \\
Dynamic & pairs only for the first 2 & 3.06 & 2.25 & 35\% & 0.9406 \\
Dynamic & compiler memories 3 of 4 & 5.54 & 4.88 & 53\% & 0.9221 \\
Dynamic & compiler memories 3 of 4 & 5.43 & 4.77 & 53\% & 0.9262 \\
Dynamic & dynamic-lookahead-layers=1 & 5.46 & 4.81 & 53\% & 0.9274 \\
Dynamic & dynamic-lookahead-layers=2 & 5.46 & 4.81 & 53\% & 0.9274 \\
Dynamic & dynamic-lookahead-layers=4 & 5.46 & 4.81 & 53\% & 0.9274 \\
Dynamic & compiler-reservation-ms=100 & 5.46 & 4.81 & 53\% & 0.9274 \\
Dynamic & compiler-reservation-ms=300 & 5.46 & 4.80 & 53\% & 0.9274 \\
ODG & baseline & 1.57 & 0.44 & 0\% & 0.9498 \\
ODG & request-duration-ms=10 & 1.57 & 0.44 & 0\% & 0.9498 \\
ODG & request-duration-ms=100 & 1.57 & 0.44 & 0\% & 0.9498 \\
ODG & request-duration-ms=20 & 1.57 & 0.44 & 0\% & 0.9498 \\
\bottomrule
\end{tabular}}
\caption{QFT, four memories, mapping seed 42, mean of three physics seeds. Each row changes one setting. These are experiments to find the cause, not proposed settings.}
\label{tab:diag-ablation}
\end{table}

Each row of Table~\ref{tab:diag-ablation} changes one setting of the QFT run at four memories.
\begin{itemize}
  \item \textbf{No effect:} the request duration (the nominal booking window, $10$ to $100$\,ms), the compiler reservation window ($100$ to $1000$\,ms) and the generation capacity per core.
    What matters is how long a booking \emph{exists}, not the length of the window it was made for. The planner's lookahead (one, two, four or eight layers) makes no difference either: not to the latency of the QFT run, and not to the planner's schedule in any of $210$ configurations (Section~\ref{sec:code-reading}, item~1).
  \item \textbf{Strong effect:} the number of memories the compiler may hold. One of four: $1.88$\,ms; two: $2.77$\,ms; three: $5.46$\,ms (ODG: $1.57$\,ms). Fewer prepared pairs, less waiting.
  \item \textbf{Preparing only first hops} (later hops fall back to on-demand): $1.89$\,ms with the compiler's three memories, and the ready pairs wait $11$ instead of $24$\,ms. This supports the chain explanation
    but does not make dynamic better than ODG at four memories; it is an experiment to isolate the cause, not a proposed policy.
\end{itemize}

\section{Other mapping seeds and other benchmarks}
\begin{table}[H]
\centering\small
\resizebox{\ifdim\width>\linewidth\linewidth\else\width\fi}{!}{%
\begin{tabular}{@{}lrrrrrr@{}}
\toprule
Policy & map 42 lat. & map 42 wait & map 43 lat. & map 43 wait & map 44 lat. & map 44 wait \\
\midrule
ODG & 1.57 & 0.44 & 1.65 & 0.52 & 1.55 & 0.42 \\
Dynamic & 5.46 & 4.81 & 5.84 & 5.19 & 5.40 & 4.73 \\
CGP & 10.15 & 9.58 & 10.26 & 9.68 & 9.94 & 9.37 \\
ACGP & 10.52 & 9.98 & 10.76 & 10.22 & 10.53 & 9.97 \\
\bottomrule
\end{tabular}}
\caption{QFT, four memories: the diagnosis on the three mapping seeds (mean of three physics seeds each; latency and wait in ms).}
\label{tab:diag-seeds}
\end{table}

\begin{table}[H]
\centering\small
\resizebox{\ifdim\width>\linewidth\linewidth\else\width\fi}{!}{%
\begin{tabular}{@{}lrrrrrrr@{}}
\toprule
Workload & ODG (ms) & Dynamic (ms) & CGP (ms) & ACGP (ms) & Dynamic hits & Dyn. stuck checks & ODG stuck checks \\
\midrule
CDKM ripple-carry adder & 1.14 & 0.30 & 0.82 & 0.76 & 95\% & 0.0\% & 0.0\% \\
VQE real amplitudes & 1.15 & 0.83 & 1.41 & 1.43 & 96\% & 4.4\% & 0.0\% \\
Random circuit & 1.11 & 0.97 & 2.08 & 2.13 & 87\% & 12.9\% & 0.0\% \\
Draper QFT adder & 1.16 & 2.45 & 5.45 & 5.65 & 73\% & 21.7\% & 0.5\% \\
Graph state & 1.19 & 4.24 & 5.03 & 5.07 & 68\% & 16.0\% & 0.7\% \\
QAOA & 1.28 & 4.07 & 8.12 & 8.43 & 56\% & 22.3\% & 3.4\% \\
QFT & 1.57 & 5.46 & 10.15 & 10.52 & 53\% & 17.2\% & 6.8\% \\
\bottomrule
\end{tabular}}
\caption{Request latency at four memories on the other benchmarks (mapping seed 42, mean of three physics seeds), the share of requests dynamic served from a ready pair, and the share of admission checks that found no free memory under dynamic and under on-demand. Dynamic is faster than on-demand where its own prepared pairs rarely get in the way.}
\label{tab:diag-workloads}
\end{table}

The picture is the same on the other two mapping seeds (Table~\ref{tab:diag-seeds}). On the other benchmarks (Table~\ref{tab:diag-workloads}) dynamic is faster than ODG on the three benchmarks where the fewest of its admission checks are stuck (the CDKM adder $0\%$, VQE $4\%$, the random circuit $13\%$)
and slower on the four where more are ($16$ to $22\%$: graph state, QFT, Draper, QAOA). That is a correlation across seven benchmarks, consistent with the mechanism, not a separate test of it. The full grid of benchmarks and memory counts is in Chapter~\ref{ch:four-policies}.

\section{What this establishes, and what it does not}
\begin{itemize}
  \item \textbf{Established} (measured, with the mechanism read in the code): at the modelled memory counts, pre-generated pairs occupy memory cards for several times as long as an on-demand request does, the
    resulting scarcity of free cards costs more waiting than the avoided generation saves, and the effect disappears when the cards are not scarce.
  \item \textbf{Established:} no tested setting of the simulator's windows or capacities explains the effect; the number of memories the compiler may hold does.
  \item \textbf{Not established:} that a better plan could not do better. The planner works in whole layers and does not know when a chain hop will actually become eligible, or that a ready pair also needs a card for its request.
    A planner that did would be a new algorithm; none is proposed or tested here.
  \item \textbf{Not established:} how this changes under other physical parameters (Chapter~\ref{ch:exp1}, sensitivity arm) or on other topologies beyond what Chapter~\ref{ch:four-policies} reports.
  \item The diagnosis rests on QFT in depth; the other benchmarks give the same sign of effect in one mapping seed.
\end{itemize}
